\chapter{کاهش‌پذیری}

در فصل قبل با قطری‌سازی ثابت کردیم $\ATM$ تصمیم‌ناپذیر است. برای مسئله‌های دیگر لازم نیست هر بار از صفر شروع کنیم. ابزار اصلی اثبات حل‌ناپذیری مسئله‌ها \term{کاهش}{Reduction} است.

\section{ایده کاهش}

\textbf{کاهش} یعنی تبدیل یک مسئله به مسئله دیگر، طوری که با داشتن راه حل مسئله دوم بتوان مسئله اول را حل کرد. اگر مسئله $A$ به مسئله $B$ کاهش یابد، می‌توانیم از راه حل $B$ برای حل $A$ استفاده کنیم. چند مثال روزمره:
\begin{itemize}
\item پیدا کردن راه در یک شهر به پیدا کردن نقشه شهر کاهش می‌یابد.
\item محاسبه مساحت مستطیل به اندازه‌گیری طول و عرض آن کاهش می‌یابد.
\item حل دستگاه معادلات خطی به وارون کردن یک ماتریس کاهش می‌یابد.
\end{itemize}

اگر $A$ به $B$ کاهش یابد، حل $A$ \textbf{سخت‌تر از} حل $B$ نیست. برای تصمیم‌پذیری این یعنی:
\begin{itemize}
\item اگر $B$ تصمیم‌پذیر باشد، $A$ هم تصمیم‌پذیر است.
\item معادلاً: اگر $A$ تصمیم‌ناپذیر باشد، $B$ هم تصمیم‌ناپذیر است.
\end{itemize}
روش کار در این فصل همیشه همین است: برای اثبات تصمیم‌ناپذیری $B$، فرض می‌کنیم $B$ تصمیم‌گیرنده‌ای دارد و با آن یک تصمیم‌گیرنده برای مسئله‌ای می‌سازیم که می‌دانیم تصمیم‌ناپذیر است (معمولاً $\ATM$). (کاهش در نظریه پیچیدگی هم نقش مهمی دارد؛ فصل ۲۴.)

\section{مسئله توقف}

\[\lng{HALT}{TM}=\{\enc{M,w}\st M\ \text{\rl{یک ماشین تورینگ است که روی ورودی $w$ متوقف می‌شود}}\}.\]

\begin{theorem}
$\lng{HALT}{TM}$ تصمیم‌ناپذیر است.
\end{theorem}

\begin{proofidea}
با برهان خلف فرض کنید ماشین $R$ زبان $\lng{HALT}{TM}$ را تصمیم بگیرد. با $R$ یک تصمیم‌گیرنده $S$ برای $\ATM$ می‌سازیم. $S$ ورودی $\enc{M,w}$ را به $R$ می‌دهد؛ $R$ همیشه متوقف می‌شود. اگر $R$ رد کند، یعنی $M$ روی $w$ متوقف نمی‌شود (حلقه می‌زند)، پس $M$ رشته $w$ را نمی‌پذیرد و $S$ رد می‌کند. اگر $R$ بپذیرد، $M$ روی $w$ متوقف می‌شود، پس می‌توانیم بدون ترس از حلقه، $M$ را روی $w$ شبیه‌سازی کنیم و همان جواب را برگردانیم.
\end{proofidea}

\begin{proof}
فرض کنید $R$ تصمیم‌گیرنده $\lng{HALT}{TM}$ باشد. ماشین $S$:
\begin{algo}{$S$}{روی ورودی $\enc{M,w}$}
\begin{enumerate}
\item $R$ را روی $\enc{M,w}$ اجرا کن.
\item اگر $R$ رد کرد، \textbf{رد کن}.
\item اگر $R$ پذیرفت، $M$ را روی $w$ شبیه‌سازی کن تا متوقف شود.
\item اگر $M$ پذیرفت \textbf{بپذیر}؛ اگر رد کرد \textbf{رد کن}.
\end{enumerate}
\end{algo}
$S$ همیشه متوقف می‌شود و $\ATM$ را تصمیم می‌گیرد، که با قضیه \ref{thm:atm} تناقض دارد. پس $\lng{HALT}{TM}$ تصمیم‌ناپذیر است.
\end{proof}

\section{تهی بودن زبان یک ماشین}

\[\lng{E}{TM}=\{\enc{M}\st M\ \text{\rl{یک ماشین تورینگ است و}}\ L(M)=\emptyset\}.\]

\begin{theorem}
$\lng{E}{TM}$ تصمیم‌ناپذیر است.
\end{theorem}

\begin{proofidea}
فرض کنید $R$ تصمیم‌گیرنده $\lng{E}{TM}$ باشد. می‌خواهیم با آن $\ATM$ را تصمیم بگیریم. روی ورودی $\enc{M,w}$، ماشین $M$ را طوری تغییر می‌دهیم که \textbf{همه رشته‌ها جز $w$ را رد کند} و روی $w$ مثل $M$ رفتار کند. زبان ماشین تغییریافته یا $\{w\}$ است (اگر $M$ رشته $w$ را بپذیرد) یا $\emptyset$. پس با پرسیدن از $R$ که «آیا زبان ماشین جدید تهی است؟» می‌فهمیم $M$ رشته $w$ را می‌پذیرد یا نه.
\end{proofidea}

\begin{proof}
فرض کنید $R$ تصمیم‌گیرنده $\lng{E}{TM}$ باشد. ماشین تغییریافته:
\begin{algo}{$M_1$}{روی ورودی $x$}
\begin{enumerate}
\item اگر $x\neq w$، \textbf{رد کن}.
\item اگر $x=w$، $M$ را روی ورودی $w$ اجرا کن و اگر $M$ پذیرفت \textbf{بپذیر}.
\end{enumerate}
\end{algo}
($M_1$ چند حالت جدید و چند بررسی برای مقایسه $x$ با $w$ دارد؛ $w$ درون توصیف $M_1$ نوشته شده.) حالا:
\begin{algo}{$S$}{روی ورودی $\enc{M,w}$}
\begin{enumerate}
\item با استفاده از توصیف $M$ و $w$، ماشین $M_1$ بالا را بساز.
\item $R$ را روی $\enc{M_1}$ اجرا کن.
\item اگر $R$ پذیرفت \textbf{رد کن}؛ اگر $R$ رد کرد \textbf{بپذیر}.
\end{enumerate}
\end{algo}
دقت کنید $S$ هرگز $M_1$ را اجرا نمی‌کند؛ فقط توصیفش را می‌سازد و به $R$ می‌دهد. اگر $M$ رشته $w$ را بپذیرد، $L(M_1)=\{w\}\neq\emptyset$ و $R$ رد می‌کند، پس $S$ می‌پذیرد. اگر نپذیرد، $L(M_1)=\emptyset$ و $S$ رد می‌کند. پس $S$ تصمیم‌گیرنده $\ATM$ است؛ تناقض.
\end{proof}

\section{آیا زبان یک ماشین منظم است؟}

\[\lng{REGULAR}{TM}=\{\enc{M}\st M\ \text{\rl{یک ماشین تورینگ است و}}\ L(M)\ \text{\rl{منظم است}}\}.\]

\begin{theorem}
$\lng{REGULAR}{TM}$ تصمیم‌ناپذیر است.
\end{theorem}
\begin{proof}
فرض کنید $R$ تصمیم‌گیرنده $\lng{REGULAR}{TM}$ باشد. روی ورودی $\enc{M,w}$، ماشین $M_2$ را می‌سازیم که اگر $M$ رشته $w$ را بپذیرد زبانی \textbf{منظم} (یعنی $\Sigma^*$) و در غیر این صورت زبانی \textbf{نامنظم} (یعنی $\{0^n1^n\st n\ge0\}$) را تشخیص دهد:
\begin{algo}{$M_2$}{روی ورودی $x$}
\begin{enumerate}
\item اگر $x$ به شکل $0^n1^n$ است، \textbf{بپذیر}.
\item در غیر این صورت $M$ را روی $w$ اجرا کن و اگر $M$ پذیرفت، \textbf{بپذیر}.
\end{enumerate}
\end{algo}
یعنی $M_2$ همه رشته‌های $\{0^n1^n\}$ را همیشه می‌پذیرد و اگر $M$ رشته $w$ را بپذیرد، بقیه رشته‌ها را هم می‌پذیرد:
\[L(M_2)=\begin{cases}\Sigma^* & \text{\rl{اگر $M$ رشته $w$ را بپذیرد}}\\ \{0^n1^n\st n\ge0\} & \text{\rl{در غیر این صورت.}}\end{cases}\]
\begin{algo}{$S$}{روی ورودی $\enc{M,w}$}
\begin{enumerate}
\item ماشین $M_2$ بالا را بساز.
\item $R$ را روی $\enc{M_2}$ اجرا کن.
\item اگر $R$ پذیرفت \textbf{بپذیر}؛ اگر رد کرد \textbf{رد کن}.
\end{enumerate}
\end{algo}
$S$ تصمیم‌گیرنده $\ATM$ است؛ تناقض.
\end{proof}

به همین شکل می‌توان ثابت کرد آزمودن اینکه زبان یک ماشین تورینگ \textbf{مستقل از متن}، \textbf{تصمیم‌پذیر} یا \textbf{متناهی} است، تصمیم‌ناپذیر است. همه این‌ها حالت‌های خاص یک قضیه کلی هستند:

\begin{theorem}[رایس]
هر ویژگی \textbf{غیربدیهی} از زبان‌هایی که ماشین‌های تورینگ تشخیص می‌دهند، تصمیم‌ناپذیر است.
\end{theorem}
در اینجا «ویژگی» یعنی مجموعه‌ای $P$ از زبان‌های تشخیص‌پذیر و می‌پرسیم آیا $L(M)\in P$؛ پس اگر $L(M_1)=L(M_2)$ آنگاه $M_1$ و $M_2$ هر دو ویژگی را دارند یا هیچ‌کدام. «غیربدیهی» یعنی دست‌کم یک ماشین ویژگی را دارد و دست‌کم یک ماشین ندارد.
\begin{proof}
فرض کنید $R$ تصمیم‌گیرنده $L_P=\{\enc{M}\st L(M)\in P\}$ باشد. بدون کاستن از کلیت فرض کنید $\emptyset\notin P$ (وگرنه همین استدلال را برای ویژگی متمم $\overline{P}$ انجام می‌دهیم؛ تصمیم‌پذیری $L_P$ و $L_{\overline P}$ معادل است). چون $P$ غیربدیهی است، ماشینی مثل $T$ با $L(T)\in P$ وجود دارد. روی ورودی $\enc{M,w}$ ماشین زیر را می‌سازیم:
\begin{algo}{$M_w$}{روی ورودی $x$}
\begin{enumerate}
\item $M$ را روی $w$ اجرا کن. اگر رد کرد \textbf{رد کن}؛ اگر پذیرفت به مرحله ۲ برو.
\item $T$ را روی $x$ اجرا کن و اگر پذیرفت \textbf{بپذیر}.
\end{enumerate}
\end{algo}
اگر $M$ رشته $w$ را بپذیرد، $L(M_w)=L(T)\in P$؛ وگرنه $L(M_w)=\emptyset\notin P$. پس با اجرای $R$ روی $\enc{M_w}$ می‌فهمیم $M$ رشته $w$ را می‌پذیرد یا نه، و $\ATM$ تصمیم‌پذیر می‌شود؛ تناقض.
\end{proof}

\section{هم‌ارزی دو ماشین}

\[\lng{EQ}{TM}=\{\enc{M_1,M_2}\st M_1,M_2\ \text{\rl{ماشین تورینگ‌اند و}}\ L(M_1)=L(M_2)\}.\]

\begin{theorem}\label{thm:eqtm}
$\lng{EQ}{TM}$ تصمیم‌ناپذیر است.
\end{theorem}
\begin{proof}
این بار از $\lng{E}{TM}$ کاهش می‌دهیم. ایده: مسئله تهی بودن حالت خاصی از هم‌ارزی است؛ کافی است ماشین را با ماشینی که همه چیز را رد می‌کند مقایسه کنیم. فرض کنید $R$ تصمیم‌گیرنده $\lng{EQ}{TM}$ باشد.
\begin{algo}{$S$}{روی ورودی $\enc{M}$}
\begin{enumerate}
\item $R$ را روی $\enc{M,M_1}$ اجرا کن که $M_1$ ماشینی است که همه ورودی‌ها را رد می‌کند.
\item اگر $R$ پذیرفت \textbf{بپذیر}؛ اگر رد کرد \textbf{رد کن}.
\end{enumerate}
\end{algo}
$S$ دقیقاً وقتی می‌پذیرد که $L(M)=L(M_1)=\emptyset$. پس $S$ تصمیم‌گیرنده $\lng{E}{TM}$ است؛ تناقض.
\end{proof}

\section{تاریخچه محاسبه و ماشین‌های خطی‌کران‌دار}

\textbf{تاریخچه محاسبه} یکی از مهم‌ترین تکنیک‌های کاهش است و معمولاً وقتی به کار می‌آید که مسئله درباره \emph{وجود} چیزی باشد (مثلاً وجود ریشه صحیح در مسئله دهم هیلبرت).

\begin{definition}[تاریخچه محاسبه]
فرض کنید $M$ ماشین تورینگ و $w$ یک ورودی باشد. یک \textbf{تاریخچه محاسبه پذیرنده} برای $M$ روی $w$ دنباله‌ای از پیکربندی‌ها $C_1,C_2,\dots,C_l$ است که $C_1$ پیکربندی شروع $M$ روی $w$، $C_l$ یک پیکربندی پذیرش و هر $C_{i+1}$ طبق قوانین $M$ از $C_i$ نتیجه می‌شود. \textbf{تاریخچه محاسبه ردکننده} هم مشابه است، با این تفاوت که $C_l$ پیکربندی رد است.
\end{definition}

تاریخچه‌های محاسبه دنباله‌های \textbf{متناهی} هستند. اگر $M$ روی $w$ متوقف نشود، هیچ تاریخچه پذیرنده یا ردکننده‌ای ندارد. ماشین‌های قطعی روی هر ورودی حداکثر یک تاریخچه محاسبه دارند؛ ماشین‌های غیرقطعی ممکن است چند تاریخچه داشته باشند (برای شاخه‌های مختلف).

\begin{definition}[ماشین خطی‌کران‌دار (\lr{LBA})]
\term{ماشین خطی‌کران‌دار}{Linear Bounded Automaton} نوعی ماشین تورینگ است که هدش اجازه ندارد از بخشی از نوار که ورودی روی آن نوشته شده بیرون برود. اگر ماشین بخواهد هد را از هر کدام از دو انتهای ورودی بیرون ببرد، هد سر جایش می‌ماند (مثل انتهای چپ نوار در ماشین تورینگ معمولی).
\end{definition}
\begin{ltrfig}
\begin{tikzpicture}[font=\small]
  \node[draw=jzMain,thick,minimum width=2cm,minimum height=0.9cm,fill=jzMain!6] (c) at (0,1.3) {control};
  \foreach \s [count=\i] in {a,b,a,b,a} {\node[cell] (t\i) at (\i*0.65,0) {$\s$};}
  \draw[->,thick,jzMain] (c.east) -| (t2.north);
\end{tikzpicture}
\end{ltrfig}

حافظه \lr{LBA} فقط به اندازه ورودی است. با بزرگ کردن الفبای نوار می‌توان حافظه را به اندازه یک \textbf{ضریب ثابت} زیاد کرد؛ پس می‌گوییم برای ورودی به طول $n$، حافظه در دسترس \textbf{خطی} بر حسب $n$ است. با این محدودیت هم \lr{LBA}ها قدرت زیادی دارند: تصمیم‌گیرنده‌های $\lng{A}{DFA}$، $\lng{A}{CFG}$، $\lng{E}{DFA}$، $\lng{E}{CFG}$ و هر زبان مستقل از متن همه \lr{LBA} هستند.

\begin{lemma}\label{lem:lba-configs}
فرض کنید $M$ یک \lr{LBA} با $q$ حالت و $g$ نماد در الفبای نوار باشد. برای نواری به طول $n$، $M$ دقیقاً $qng^n$ پیکربندی متمایز دارد.
\end{lemma}
\begin{proof}
پیکربندی تصویری از ماشین در میانه محاسبه است و از حالت کنترل، مکان هد و محتوای نوار تشکیل می‌شود. $M$ دارای $q$ حالت است، هد می‌تواند در یکی از $n$ خانه باشد و $g^n$ رشته ممکن از نمادهای نوار می‌تواند روی نوار باشد. حاصل‌ضرب این سه عدد تعداد کل پیکربندی‌هاست.
\end{proof}

\begin{theorem}
$\lng{A}{LBA}=\{\enc{M,w}\st M\ \text{\rl{یک \lr{LBA} است که $w$ را می‌پذیرد}}\}$ \textbf{تصمیم‌پذیر} است.
\end{theorem}
\begin{proofidea}
اگر $M$ یک پیکربندی را تکرار کند، از آن به بعد همان مسیر را دوباره و دوباره طی می‌کند، یعنی در حلقه افتاده است. چون $M$ فقط تعداد محدودی پیکربندی دارد، اگر بیش از آن تعداد گام اجرا شود حتماً پیکربندی تکراری داشته است. پس کافی است $M$ را حداکثر $qng^n$ گام اجرا کنیم.
\end{proofidea}
\begin{proof}
\begin{algo}{$L$}{روی ورودی $\enc{M,w}$ که $M$ یک \lr{LBA} است}
\begin{enumerate}
\item $M$ را روی $w$ به اندازه $qng^n$ گام یا تا وقتی متوقف شود شبیه‌سازی کن ($n=\abs{w}$).
\item اگر $M$ متوقف شد، اگر پذیرفته بود \textbf{بپذیر} و اگر رد کرده بود \textbf{رد کن}. اگر متوقف نشد، \textbf{رد کن}.
\end{enumerate}
\end{algo}
اگر $M$ در $qng^n$ گام متوقف نشود، طبق لم \ref{lem:lba-configs} و اصل لانه کبوتری یک پیکربندی را تکرار کرده و هرگز متوقف نمی‌شود؛ پس $w$ را نمی‌پذیرد و رد کردن درست است.
\end{proof}

\begin{theorem}
$\lng{E}{LBA}=\{\enc{M}\st M\ \text{\rl{یک \lr{LBA} است و}}\ L(M)=\emptyset\}$ تصمیم‌ناپذیر است.
\end{theorem}

\begin{proofidea}
از $\ATM$ کاهش می‌دهیم. روی ورودی $\enc{M,w}$ یک \lr{LBA} مثل $B$ می‌سازیم که دقیقاً \textbf{تاریخچه‌های محاسبه پذیرنده $M$ روی $w$} را می‌پذیرد. اگر $M$ رشته $w$ را بپذیرد، زبان $B$ یک رشته دارد (چون $M$ قطعی است)؛ وگرنه زبان $B$ تهی است. پس با پرسیدن تهی بودن زبان $B$ می‌فهمیم $M$ رشته $w$ را می‌پذیرد یا نه.

$B$ رشته‌ای به شکل زیر را دریافت می‌کند:
\[\#\,C_1\,\#\,C_2\,\#\,C_3\,\#\cdots\#\,C_l\,\#\]
و سه شرط را بررسی می‌کند:
\begin{enumerate}
\item $C_1$ پیکربندی شروع $M$ روی $w$ است.
\item هر $C_{i+1}$ طبق قوانین $M$ از $C_i$ نتیجه می‌شود.
\item $C_l$ یک پیکربندی پذیرش $M$ است.
\end{enumerate}
\end{proofidea}

\begin{proof}
\textbf{ساختن $B$:} بررسی شرط ۱: $B$ رشته $C_1$ را با $q_0w_1w_2\cdots w_n$ مقایسه می‌کند ($w$ در توصیف $B$ نوشته شده). شرط ۳: $B$ بررسی می‌کند که $C_l$ حالت $\qacc$ را دارد. شرط ۲: برای هر $i$، $B$ بین $C_i$ و $C_{i+1}$ رفت و برگشت می‌کند و نمادهای متناظر را با علامت‌گذاری روی نوار (که همان ورودی است) مقایسه می‌کند. دو پیکربندی متوالی باید جز در اطراف هد یکسان باشند و تغییر اطراف هد باید طبق تابع گذار $M$ باشد. همه این کارها بدون بیرون رفتن از بخش ورودی نوار انجام می‌شود، پس $B$ یک \lr{LBA} است.

\textbf{کاهش:} فرض کنید $R$ تصمیم‌گیرنده $\lng{E}{LBA}$ باشد.
\begin{algo}{$S$}{روی ورودی $\enc{M,w}$}
\begin{enumerate}
\item \lr{LBA} $B$ را از روی $M$ و $w$ بساز.
\item $R$ را روی $\enc{B}$ اجرا کن.
\item اگر $R$ رد کرد \textbf{بپذیر}؛ اگر پذیرفت \textbf{رد کن}.
\end{enumerate}
\end{algo}
اگر $M$ رشته $w$ را بپذیرد، تاریخچه پذیرنده‌ای وجود دارد، $L(B)\neq\emptyset$ و $S$ می‌پذیرد؛ وگرنه $L(B)=\emptyset$ و $S$ رد می‌کند. پس $S$ تصمیم‌گیرنده $\ATM$ است؛ تناقض.
\end{proof}

\section{گرامرهایی که همه رشته‌ها را تولید می‌کنند}

\[\lng{ALL}{CFG}=\{\enc{G}\st G\ \text{\rl{یک \lr{CFG} است و}}\ L(G)=\Sigma^*\}.\]

\begin{theorem}
$\lng{ALL}{CFG}$ تصمیم‌ناپذیر است.
\end{theorem}

\begin{proofidea}
باز هم از $\ATM$ با تاریخچه محاسبه کاهش می‌دهیم. روی ورودی $\enc{M,w}$ یک گرامر $G$ می‌سازیم که:
\begin{itemize}
\item اگر $M$ رشته $w$ را نپذیرد، $G$ \textbf{همه} رشته‌ها را تولید کند؛
\item اگر $M$ رشته $w$ را بپذیرد، $G$ همه رشته‌ها \textbf{جز} تاریخچه محاسبه پذیرنده $M$ روی $w$ را تولید کند.
\end{itemize}
پس $L(G)=\Sigma^*$ اگر و تنها اگر $M$ رشته $w$ را نپذیرد.
\end{proofidea}

\begin{proof}
یک تاریخچه محاسبه پذیرنده را به شکل $\#C_1\#C_2\#\cdots\#C_l\#$ می‌نویسیم. $G$ همه رشته‌هایی را تولید می‌کند که دست‌کم یکی از این شرط‌ها را داشته باشند:
\begin{enumerate}
\item با $C_1$ (پیکربندی شروع $M$ روی $w$) شروع نشوند،
\item به یک پیکربندی پذیرش ختم نشوند، یا
\item در آن‌ها یک $C_i$ وجود داشته باشد که $C_{i+1}$ به‌درستی از آن طبق قوانین $M$ به دست نیامده باشد.
\end{enumerate}
اگر $M$ رشته $w$ را نپذیرد، هیچ تاریخچه پذیرنده‌ای وجود ندارد، پس \emph{هر} رشته دست‌کم یکی از این شرط‌ها را دارد و $G$ همه رشته‌ها را تولید می‌کند. اگر $M$ رشته $w$ را بپذیرد، تنها رشته‌ای که هیچ‌کدام از شرط‌ها را ندارد همان تاریخچه پذیرنده است.

ساختن مستقیم $G$ دشوار است. به جای آن یک \lr{PDA} مثل $D$ می‌سازیم و آن را با قضیه \ref{thm:pda-cfg} به گرامر تبدیل می‌کنیم. $D$ به‌طور غیرقطعی یکی از سه شرط را انتخاب و بررسی می‌کند:
\begin{enumerate}
\item ابتدای ورودی را بررسی می‌کند و اگر $C_1$ پیکربندی شروع درست نبود، می‌پذیرد.
\item انتهای ورودی را بررسی می‌کند و اگر آخرین پیکربندی، پیکربندی پذیرش نبود، می‌پذیرد.
\item به‌طور غیرقطعی یک $C_i$ را انتخاب می‌کند، آن را روی پشته می‌گذارد و هنگام خواندن $C_{i+1}$ آن را از پشته برمی‌دارد و مقایسه می‌کند. اگر ناهماهنگی یا به‌روزرسانی نادرستی (طبق تابع گذار $M$ در اطراف هد) پیدا کرد، می‌پذیرد.
\end{enumerate}
یک مشکل فنی در شرط ۳: وقتی $C_i$ را از پشته برمی‌داریم، به ترتیب معکوس بیرون می‌آید. برای حل این مشکل تاریخچه را طوری می‌نویسیم که پیکربندی‌های زوج \textbf{معکوس} باشند:
\[\#\,\overrightarrow{C_1}\,\#\,\overleftarrow{C_2^R}\,\#\,\overrightarrow{C_3}\,\#\,\overleftarrow{C_4^R}\,\#\cdots\#\,C_l\,\#\]
به این ترتیب هر دو پیکربندی متوالی را می‌توان با پشته مقایسه کرد.

حالا اگر $R$ تصمیم‌گیرنده $\lng{ALL}{CFG}$ باشد، ماشین $S$ روی $\enc{M,w}$ گرامر $G$ را از روی $D$ می‌سازد، $R$ را روی $\enc{G}$ اجرا می‌کند و اگر $R$ رد کرد می‌پذیرد و اگر پذیرفت رد می‌کند. $S$ تصمیم‌گیرنده $\ATM$ است؛ تناقض.
\end{proof}

\begin{corollary}
$\lng{EQ}{CFG}$ تصمیم‌ناپذیر است.
\end{corollary}
\begin{proof}
اگر تصمیم‌گیرنده‌ای برای $\lng{EQ}{CFG}$ داشتیم، می‌توانستیم $\lng{ALL}{CFG}$ را تصمیم بگیریم: روی ورودی $\enc{G}$، گرامر $G$ را با گرامر ساده $H:\ S\to aS\mid\eps$ (برای هر $a\in\Sigma$) که $\Sigma^*$ را تولید می‌کند مقایسه می‌کنیم. $\enc{G,H}\in\lng{EQ}{CFG}$ اگر و تنها اگر $L(G)=\Sigma^*$.
\end{proof}

\begin{summary}
\begin{itemize}
\item اگر $A$ به $B$ کاهش یابد و $A$ تصمیم‌ناپذیر باشد، $B$ هم تصمیم‌ناپذیر است.
\item تصمیم‌ناپذیر (با کاهش از $\ATM$): $\lng{HALT}{TM}$، $\lng{E}{TM}$، $\lng{REGULAR}{TM}$ و هر ویژگی غیربدیهی زبان ماشین‌ها (رایس). $\lng{EQ}{TM}$ با کاهش از $\lng{E}{TM}$.
\item \lr{LBA} تعداد متناهی ($qng^n$) پیکربندی دارد، پس $\lng{A}{LBA}$ تصمیم‌پذیر است؛ اما $\lng{E}{LBA}$ تصمیم‌ناپذیر است.
\item با تاریخچه محاسبه: $\lng{E}{LBA}$ و $\lng{ALL}{CFG}$ (و در نتیجه $\lng{EQ}{CFG}$) تصمیم‌ناپذیرند.
\end{itemize}
\end{summary}
